\documentclass[11pt]{article}
\usepackage[margin=1in]{geometry}
\usepackage{amsmath,amssymb,amsthm}
\usepackage[T1]{fontenc}
\usepackage{lmodern}
\usepackage[hidelinks]{hyperref}
\newtheorem{theorem}{Theorem}
\newtheorem{lemma}[theorem]{Lemma}
\newcommand{\R}{\mathbb R}
\newcommand{\Z}{\mathbb Z}
\title{Exact chordal separation of spherical Fibonacci points}
\author{Alec Kriebel\\Independent researcher\\
\small\url{https://orcid.org/0009-0001-9320-500X}}
\date{October 5, 2026}
\begin{document}
\maketitle
\begin{abstract}
We prove the exact finite separation conjecture stated in Brauchart's
2012 Oberwolfach contribution: the rational, unshifted spherical Fibonacci
configuration with $F_n$ points has minimum Euclidean chord $2/\sqrt{F_n}$
for every $n\geq3$. A latitude--longitude inequality reduces the lower
bound to the classical product separation of Fibonacci lattices. We
include an elementary proof of the weaker arithmetic estimate needed
here, making the argument self-contained.
\end{abstract}

\section{The conjecture and its precise scope}
Let $F_0=0$, $F_1=F_2=1$, and $F_{n+1}=F_n+F_{n-1}$. For an integer
$n\geq3$, put $q=F_n$ and $p=F_{n-1}$, and define
\begin{equation}\label{eq:points}
 z_k=\left(2\sqrt{\frac{k}{q}\left(1-\frac{k}{q}\right)}
 \cos\frac{2\pi pk}{q},\quad
 2\sqrt{\frac{k}{q}\left(1-\frac{k}{q}\right)}
 \sin\frac{2\pi pk}{q},\quad 1-\frac{2k}{q}\right),
 \qquad 0\leq k<q.
\end{equation}
These are distinct points on the unit sphere. Brauchart's contribution
\cite[pp.~2436--2439, especially p.~2438]{Brauchart} conjectures the
following exact minimum, already attained by the north pole and its
first latitude point.
\begin{theorem}\label{thm:main}
For every integer $n\geq3$, the configuration \eqref{eq:points} satisfies
\[
 \min_{0\leq i<j<q}\|z_i-z_j\|=\frac{2}{\sqrt q}.
\]
The unique minimizing unordered pair is $\{z_0,z_1\}$.
\end{theorem}
Here the metric is the ambient Euclidean chord. Its corresponding
minimum geodesic distance is $2\arcsin(1/\sqrt q)$. The singleton $n=2$
has no pair and is excluded. The theorem concerns the displayed rational
longitudes and unshifted heights; it makes no assertion about midpoint
heights, irrational golden angles, arbitrary cardinalities, discrepancy,
or optimal packing among all spherical configurations.

\section{The arithmetic and geometric inequalities}
The arithmetic ingredient is classical. For the dual Fibonacci lattice
$\{(h_1,h_2)\in\Z^2:h_1+ph_2\equiv0\pmod q\}$, its hyperbolic-cross
index, the minimum of
$\max(1,|h_1|)\max(1,|h_2|)$ over nonzero vectors, is $F_{n-2}$;
see Zaremba \cite[p.~186]{Zaremba} and the explicit modern statement in
Suryanarayana, Cools and Nuyens \cite[Theorem~2.2]{SCN}.
Applying this to $(-r,m)$ gives $m|r|\geq F_{n-2}\geq q/3$.
For completeness, the following proof establishes the needed weaker
inequality without relying on that index theorem.

\begin{lemma}\label{lem:arithmetic}
If $n\geq3$, $1\leq m<q=F_n$, and $\ell\in\Z$, then
\[
 m\,|mp-\ell q|\geq q/3,\qquad p=F_{n-1}.
\]
\end{lemma}
\begin{proof}
Choose the unique $j\geq2$ with $F_j\leq m<F_{j+1}$.
Then $j\leq n-1$; write $h=n-j\geq1$.
The integer vectors $(F_j,F_{j-1})$ and $(F_{j+1},F_j)$ form a
unimodular basis, since their determinant is $(-1)^{j-1}$.
Consequently
\[
 (m,\ell)=S(F_j,F_{j-1})+T(F_{j+1},F_j),\qquad S,T\in\Z.
\]
Here $S\ne0$: otherwise a positive $m<F_{j+1}$ would be a multiple
of $F_{j+1}$. Both $S,T$ positive would give $m\geq F_j+F_{j+1}$;
both negative, or $S<0,T=0$, would give $m<0$. Thus either $T=0,S>0$
or $S,T$ have opposite signs. The recurrence gives
\[
 pF_j-qF_{j-1}=(-1)^{j-1}F_h,\qquad
 pF_{j+1}-qF_j=(-1)^jF_{h-1}.
\]
It follows that $|mp-\ell q|=|SF_h-TF_{h-1}|\geq F_h$,
including the terminal case $h=1$, where $F_{h-1}=0$.
Finally the Fibonacci addition identity gives
\[
 q=F_jF_{h+1}+F_{j-1}F_h\leq3F_jF_h
 \leq3m|mp-\ell q|,
\]
using $F_{h+1}\leq2F_h$ and $F_{j-1}\leq F_j$.
The determinant, error and addition identities all follow by induction
from the defining recurrence.
\end{proof}

\begin{lemma}\label{lem:geometry}
For two points with latitude parameters $0\leq a<b\leq1$ in
\eqref{eq:points}, let $d=b-a$ and let $c$ be the cosine of their
longitude difference. Their squared chord $D^2$ satisfies
\[
 c\leq0\ \Longrightarrow\ D^2\geq4d,
 \qquad
 0<c<1\ \Longrightarrow\ D^2\geq4d\sqrt{1-c^2}.
\]
\end{lemma}
\begin{proof}
Put $U=2(a+b-2ab)$ and $V=4\sqrt{a(1-a)b(1-b)}$.
Direct expansion gives
\[
 D^2=2(U-cV),\quad U\geq2d,\quad V^2=U^2-4d^2,
 \quad U-2d=4a(1-b).
\]
For $c\leq0$ the first bound follows. For $0<c<1$, $U-cV>0$,
and the identity
\[
 (U-cV)^2-4d^2(1-c^2)=(V-cU)^2\geq0
\]
proves the second bound. No planar distance-preservation property of
the equal-area spherical map is used.
\end{proof}

\section{Proof of the exact separation}
\begin{proof}[Proof of Theorem~\ref{thm:main}]
Fix $0\leq i<j<q$ and set $m=j-i$ and $d=m/q$.
Choose $\ell\in\Z$ so that $r=mp-\ell q$ satisfies
$0<|r|\leq q/2$. Such a nonzero residue exists because consecutive
Fibonacci numbers are coprime and $0<m<q$.
The longitude cosine is $c=\cos(2\pi r/q)$.
If $c\leq0$, Lemma~\ref{lem:geometry} gives
$\|z_i-z_j\|^2\geq4m/q\geq4/q$.
If $c>0$, then $0<|r|<q/4$. Concavity of sine on $[0,\pi/2]$ gives
$\sin(2\pi|r|/q)\geq4|r|/q$, so the two lemmas give
\[
 \|z_i-z_j\|^2\geq\frac{4m}{q}\sin\frac{2\pi|r|}{q}
 \geq\frac{16m|r|}{q^2}\geq\frac{16}{3q}>\frac4q.
\]
This proof includes $q=2,3,5$ without separate numerical cases.
Since $z_0=(0,0,1)$, direct calculation gives
$\|z_0-z_1\|^2=4/q$, proving the exact minimum.
For equality in the nonpositive-cosine case we must have $m=1$ and
$U=2d$. As $b=j/q<1$, the identity $U-2d=4a(1-b)$ forces $a=0$.
Thus equality occurs only for $i=0,j=1$.
\end{proof}

\section{Source conventions, priority and verification}
The displayed range $0\leq k<q$ and explicit pole witness in
\cite[p.~2438]{Brauchart} specify the convention used here, despite
the following set notation's endpoint shift. The range $1\leq k\leq q$
gives an isometric set via $k\mapsto q-k$ and
$(x,y,z)\mapsto(x,-y,-z)$. Including both endpoints instead gives $q+1$
points. The reversed planar coordinate order in
\cite[Section~5.2]{ABD} gives an isometric unordered set: reindex
$k\equiv pj\pmod q$ and use $p^2\equiv(-1)^n\pmod q$; odd $n$ also
requires reflection of the second spherical coordinate.

The construction, pole upper bound and arithmetic product estimate are
prior work. The contribution here is the all-pair spherical lower bound
and its proof of the precisely stated conjecture. An extensive bounded
primary-source audit found no earlier exact resolution in the sources
inspected. It does not establish worldwide first priority or certify
the conjecture's openness as of 2026. The fate of the Brauchart--Dick
in-preparation item cited in \cite{Brauchart}, and some final graphics
editions, remains unverified; the accompanying priority note records
the actual reading and access limits.

The accompanying verification package supplies exact finite arithmetic
and certified rational chord controls, conventions and counterexample
controls, source provenance, and reproduction instructions. These are
checks supplementary to the proof, not a computational proof of an
infinite theorem. AI tools were used extensively for derivation,
literature comparison, computation, writing and independent adversarial
verification. This is an unrefereed preprint; no human peer review or
formal proof-assistant certification is claimed.

\begin{thebibliography}{9}
\bibitem{Brauchart} J.~S.~Brauchart,
Spherical Nets, Sequences and Lattice rules---construction of
low-discrepancy sequences on spheres, in \emph{Optimal and Near Optimal
Configurations on Lattices and Manifolds}, Oberwolfach Reports
\textbf{9} (2012), 2436--2439; published 2013.
\href{https://doi.org/10.4171/OWR/2012/40}{doi:10.4171/OWR/2012/40}.
\bibitem{ABD} C.~Aistleitner, J.~S.~Brauchart and J.~Dick,
Point sets on the sphere $\mathbb S^2$ with small spherical cap
discrepancy, \emph{Discrete Comput.~Geom.} \textbf{48} (2012), 990--1024.
\href{https://doi.org/10.1007/s00454-012-9451-3}{doi:10.1007/s00454-012-9451-3}.
\bibitem{Zaremba} S.~K.~Zaremba,
A remarkable lattice generated by Fibonacci numbers,
\emph{Fibonacci Quart.} \textbf{8} (1970), 185--198.
\url{https://www.fq.math.ca/Scanned/8-2/zaremba.pdf}.
\bibitem{SCN} G.~Suryanarayana, R.~Cools and D.~Nuyens,
Integration and Approximation with Fibonacci lattice points,
\emph{Dolomites Res.~Notes Approx.} \textbf{8} (2015), 92--101.
\url{https://www.maths.tcd.ie/EMIS/journals/DRNA/pdf/cools_etal_10ypdpts.pdf}.
\end{thebibliography}
\end{document}
